% !TeX root = thesis.tex

\chapter{Conclusions and Future Work} \label{chap:conclusions}

\section{Conclusions} \label{sec:conclusions}

This dissertation set out to develop and evaluate a multi-task DL framework capable of jointly
performing AD, FD and PdM for the vacuum pump of a
biomethane upgrading plant, and to compare it against single-task DL models and traditional
ML baselines. The research gap identified in Chapter~\ref{chap:sota} was twofold: existing
multi-task approaches to industrial condition monitoring combine at most two of these three tasks, and
none of the approaches reviewed had been applied to a biomethane upgrading context, where no prior work
addresses the vacuum pump specifically.

Addressing this required first constructing a dataset that did not exist. One month of real operational
data was obtained from a VPSA plant in Lithuania, but because no vacuum pump fault has ever been recorded
across the operator's monitored plants, the data contains only normal operation. A cleaning pipeline was
developed that partitions the record into eleven independent production sessions, and a fault injection
framework was built on top of it that generates four fault scenarios at three severity levels, two
combined scenarios in which faults overlap in time, and a false positive scenario, all calibrated to the
warning and trip thresholds specified by the pump manufacturer. Together with the frozen sensor episode
identified during exploratory analysis and deliberately retained, this yields a labelled dataset of
145{,}826 observations supporting all three tasks simultaneously.

Five models were then trained and evaluated on this dataset across eleven window sizes: two traditional
baselines, three single-task LSTM models, and two multi-task LSTM architectures differing in how the three
task-specific heads are organised around a shared backbone.

Three conclusions follow from the results. First, the recurrent models substantially outperformed the traditional baselines on all three tasks. The
margin was largest on PdM.

Second, and central to the research question, a single multi-task model performed comparably to three
independently trained single-task models. The single-task models obtained the best anomaly detection
F1-score and the best PdM $R^2$, while the cascaded multi-task
architecture obtained the best FD mean F1-score. Since the multi-task architecture covers all three tasks with one shared backbone and a
single forward pass, rather than three separate models to train, tune and maintain, this equivalence
favours the multi-task formulation for a deployed monitoring system.

Third, of the two multi-task architectures, the cascaded design matched or exceeded the joint design on
all three tasks, with the largest difference on FD. Conditioning each head on the predictions
of the preceding ones, mirroring the diagnostic sequence of first establishing whether something is wrong,
then what it is, and only then how severe it is, therefore proved to be a better way of organising the
three heads than branching them independently off the shared representation.

Alongside these findings, the results also identify where the approach is weakest. The bearing fault, the
rarest and shortest of the four scenarios, was detected poorly by every model, and worse by the multi-task
architectures than by the single-task model or the baselines. Since the conclusions rest on synthetically injected faults
evaluated on two sessions of a single plant, they establish that the approach is viable rather than that
it is validated in operation.

\section{Future Work} \label{sec:future_work}

\subsection*{Deployment on the plant}

The natural continuation of this work is to deploy the trained models on live plant data. This has not yet
been carried out because the operator's server is currently undergoing an infrastructure upgrade and is
not operational for this purpose, but deployment is planned once the upgrade is complete. Running the
models against the live SCADA stream would allow their behaviour to be observed under real operating
conditions and, in particular, would allow the false positive rate to be measured over an extended period
of normal operation, which is the property that most directly determines whether operators come to trust
such a system. The models developed here are compatible with this setting: they consume the same fourteen
process variables the plant already records at one-minute resolution, and produce all three outputs in a
single forward pass.

\subsection*{Validation against real faults}

The most significant limitation of this dissertation is the absence of real fault data. As failures are
eventually recorded on the monitored plants, the injected scenarios should be compared against the
observed signatures, so that the fault injection framework can be recalibrated where it diverges and the
models retrained on genuine occurrences. This would also allow the performance reported here to be
reassessed against real degradation rather than against designed deviations.

\subsection*{Extending the models}

The hyperparameter search was intentionally small, and its largest configuration was frequently selected
as the best, suggesting that larger backbones may yield further gains. Beyond simply enlarging the search,
architectures better suited to short transient events would be worth investigating, given that the bearing
fault proved difficult for every model: a convolutional stage preceding the recurrent layers, or an
attention mechanism over the input window, would both provide a more direct means of localising a brief
event within a long sequence. The fixed loss weights used by the multi-task models could also be replaced
by an adaptive weighting scheme, which may reduce the penalty the rarest fault type currently incurs under
the shared representation.

\subsection*{Broadening the scope}

Finally, the framework is not specific to the vacuum pump. The same pipeline, consisting of session-based
cleaning, threshold-calibrated fault injection, and a multi-task recurrent model, could be applied to other
critical assets in the VPSA process, such as the compressors or the adsorption valves, and to the plants
the operator monitors elsewhere. Applying it across multiple plants would additionally allow the question
this dataset cannot answer to be addressed, namely whether a model trained on one installation transfers
to another.